\chapter{Deep Q-networks}
\label{ch:dqn}

Tabular methods store one value per state; Pacman has about $10^{30}$ of
them (\cref{pac:n-states}). The way out is to represent $q$ by a
parametrised function $Q_\theta$, here a convolutional network reading the
image of the board, and to adjust the parameters $\theta$ so that
$Q_\theta$ approximately satisfies the Bellman optimality equation. States
that look alike then share what is learnt about them. This chapter shows why
the naive combination of Q-learning and neural networks is unstable, how the
deep Q-network (DQN) of Mnih et al.\ (2015) stabilises it, and the
improvements used by the DQN agent of the repository.

\section{Q-learning with function approximation}

With a parametrised $Q_\theta$, the update \eqref{eq:q-learning} becomes a
gradient step on the squared TD error, the target being treated as a
constant:
\begin{equation}
  \label{eq:semi-gradient}
  \theta \gets \theta + \alpha \left[ R_{t+1} + \gamma \max_{a'}
    Q_\theta(S_{t+1}, a') - Q_\theta(S_t, A_t) \right]
    \nabla_\theta Q_\theta(S_t, A_t).
\end{equation}
This is not the gradient of any fixed objective (the target also depends on
$\theta$, and its dependence is ignored): it is called a
\emph{semi-gradient} method. With a table, $\nabla_\theta Q_\theta(s, a)$
is the indicator of the entry $(s, a)$, and \eqref{eq:semi-gradient} is
exactly tabular Q-learning. With a network, changing $\theta$ to fit one
state changes the values of all the others, including the target.

\subsection{The deadly triad}

Sutton and Barto call the combination of \emph{function approximation},
\emph{bootstrapping} and \emph{off-policy} learning the \emph{deadly
triad}: each is harmless alone, together they can make the parameters
diverge. A two-state example shows the mechanism.

\begin{proposition}{Divergence of semi-gradient TD}{divergence}
  Two states $s_1, s_2$ have features $x(s_1) = 1$ and $x(s_2) = 2$, and
  values $V_w(s) = w\, x(s)$ with a single parameter $w$. The transition
  $s_1 \to s_2$ gives reward $0$. If the semi-gradient TD(0) update is
  applied repeatedly to this transition only (an off-policy distribution of
  updates), then
  \begin{equation*}
    w_{k+1} = \left(1 + \alpha (2 \gamma - 1)\right) w_k,
  \end{equation*}
  and $\abs{w_k} \to \infty$ for every $w_0 \neq 0$ as soon as $\gamma >
  1/2$ and $\alpha > 0$.
\end{proposition}

\begin{proof}
  The update is $w \gets w + \alpha [0 + \gamma\, V_w(s_2) - V_w(s_1)]\,
  \nabla_w V_w(s_1) = w + \alpha (2 \gamma w - w) \cdot 1$. The factor $1 +
  \alpha (2 \gamma - 1)$ exceeds $1$ when $\gamma > 1/2$.
\end{proof}

Each update tries to raise $V(s_1)$ towards $\gamma V(s_2)$; since $s_2$
shares the parameter and has a larger feature, its value rises even more,
and so does the target. The transition from $s_2$, which would correct
this, is never used. DQN does not remove the deadly triad, but its two key
ingredients make such feedback loops much weaker.

\section{DQN}

\begin{definition}{Experience replay and target network}{dqn}
  \begin{itemize}[nosep]
    \item The \emph{replay memory} stores the last transitions $(s, a, r,
      s', d)$ played ($d = 1$ if $s'$ is terminal). Each update uses a
      mini-batch sampled uniformly from it.
    \item The \emph{target network} $Q_{\theta^-}$ is a copy of $Q_\theta$
      used to compute the targets; it is updated slowly (every $C$ steps,
      or by Polyak averaging, see below).
  \end{itemize}
  The update minimises over the mini-batch the loss
  \begin{equation}
    \label{eq:dqn-loss}
    \mathcal{L}(\theta) = \frac{1}{B} \sum_{i=1}^{B} \ell\left(y_i -
      Q_\theta(s_i, a_i)\right), \qquad
    y_i = r_i + \gamma (1 - d_i) \max_{a'} Q_{\theta^-}(s'_i, a'),
  \end{equation}
  with the targets $y_i$ held constant, and $\ell$ the squared error or
  the Huber loss.
\end{definition}

\paragraph{Why a replay memory.} Consecutive transitions of a game are
strongly correlated: a mini-batch made of them describes one situation and
pushes the network towards it, then the next situation pushes it
elsewhere. Sampling from a large memory gives nearly independent samples
covering many situations, closer to the assumptions of stochastic gradient
descent. It also lets each transition be used several times, which matters
when playing is expensive. This is only possible because Q-learning is
off-policy: the transitions were generated by older policies, but the
target is the Bellman optimality target, valid whatever the policy that
produced the data.

\paragraph{Why a target network.} With the online network in the target,
each gradient step moves the target (\cref{prop:divergence}). With a frozen
copy, each phase between two copies is an ordinary regression problem
towards fixed targets, approximately one application of $\T^\star$:
DQN behaves like an approximate value iteration.

\paragraph{The Huber loss.} The Huber loss $\ell(u) = u^2 / 2$ for $\abs{u}
\leq 1$ and $\abs{u} - 1/2$ otherwise has derivative $\ell'(u) = \max(-1,
\min(1, u))$: errors larger than $1$ contribute a bounded gradient. Large TD
errors, frequent at the beginning of training or after a rare event (a
death worth $-35$), then do not destabilise the network.

\begin{algobox}{DQN (with the improvements of \cref{sec:dqn-improvements})}
  \begin{algorithmic}[1]
    \State initialise $\theta$; $\theta^- \gets \theta$; empty replay memory
    \For{each step $t$}
      \State play $a_t$, $\varepsilon$-greedy with respect to
        $Q_\theta(s_t, \cdot)$; observe $r_{t+1}, s_{t+1}$
      \State store the transition (as an $n$-step transition) in the memory
      \State sample a mini-batch; compute the targets with $Q_{\theta^-}$
        (double DQN target \eqref{eq:double-target})
      \State one gradient step on \eqref{eq:dqn-loss} (Huber loss, gradient
        clipping)
      \State $\theta^- \gets \tau \theta + (1 - \tau) \theta^-$
      \State decrease $\varepsilon$ (linear schedule)
    \EndFor
  \end{algorithmic}
\end{algobox}

\section{Improvements}
\label{sec:dqn-improvements}

\subsection{Double DQN}

The target \eqref{eq:dqn-loss} takes a maximum of noisy estimates, which
overestimates (\cref{prop:max-bias}). Double DQN (van Hasselt et al., 2016)
uses the online network to \emph{choose} the next action and the target
network to \emph{evaluate} it:
\begin{equation}
  \label{eq:double-target}
  y = r + \gamma (1 - d)\, Q_{\theta^-}\!\left(s', \argmax_{a'}
    Q_\theta(s', a')\right).
\end{equation}
When the errors of the two networks are not too correlated, a noisy
overestimate by the online network is not confirmed by the target network.

\subsection{Dueling architecture}

In many states the choice of action barely matters (in an empty corridor
far from the ghosts, going on or back changes little), while the value of
the state matters a lot. The \emph{dueling} network (Wang et al., 2016)
computes a state value $V(s)$ and advantages $A(s, a)$ with two heads on a
common trunk, and combines them as
\begin{equation}
  \label{eq:dueling}
  Q(s, a) = V(s) + A(s, a) - \frac{1}{\abs{\Actions}} \sum_{a'} A(s, a').
\end{equation}

\begin{proposition}{Identifiability of the dueling decomposition}{dueling}
  For every function $Q$, the decomposition \eqref{eq:dueling} with the
  constraint $\sum_a A(s, a) = 0$ exists and is unique: $V(s) =
  \frac{1}{\abs{\Actions}} \sum_a Q(s, a)$ and $A(s, a) = Q(s, a) - V(s)$.
  Without the mean subtraction, $(V + c, A - c)$ gives the same $Q$ for
  every constant $c$.
\end{proposition}

\begin{proof}
  Averaging \eqref{eq:dueling} over $a$ gives $\frac{1}{\abs{\Actions}}
  \sum_a Q(s, a) = V(s)$: the advantage terms cancel. Then $A(s, a) - \bar
  A(s) = Q(s, a) - V(s)$, where $\bar A(s)$ is the mean of the advantages;
  with the constraint $\bar A(s) = 0$ this determines $A$. The last claim
  is immediate.
\end{proof}

Every transition updates $V(s)$, whatever the action played, so the value
of a state is learnt from all the actions taken in it.

\subsection{$n$-step targets}

The repository's DQN uses $n$-step targets (\cref{prop:nstep-error}): for a
transition from $s_t$,
\begin{equation}
  y = \sum_{k=0}^{m-1} \gamma^k r_{t+k+1} + \gamma^m (1 - d)\,
  Q_{\theta^-}\!\left(s_{t+m}, \argmax_{a'} Q_\theta(s_{t+m}, a')\right),
\end{equation}
with $m = n = 3$, or $m < n$ if the game ended (or was aborted) before. The
replay memory assembles these transitions on the fly. The $n$-step target
is not exactly off-policy any more (the intermediate actions were chosen by
an older, exploratory policy); with small $n$ this bias is outweighed by
the faster propagation of rewards.

The replay memory keeps, for each game, the last $n$ steps not yet
stored. When it has $n$ of them, or when the game ends (or is aborted), it
stores the oldest one as an $n$-step transition: discounted sum of the
rewards, state reached, \code{done} flag, and the discount $\gamma^m$ to
apply to the bootstrap.
\begin{lstlisting}[caption={Assembling $n$-step transitions
  (\code{memories.py}, \code{NStepReplayBuffer})}]
def _emit_first(self, env_idx, s_next, done):
    q = self.pending[env_idx]
    ret = 0.
    for k, (_, _, r) in enumerate(q):     # sum_k gamma^k r_{t+k+1}
        ret += (self.gamma ** k) * r
    s, a, _ = q.popleft()
    self._store(s, a, ret, s_next, done, self.gamma ** (len(q) + 1))

def push(self, env_idx, state, action, reward, next_state, done,
         truncated=False):
    q = self.pending[env_idx]
    q.append((state, action, reward))
    if done or truncated:     # flush: shorter horizons; done only if
        while q:              # the game really ended (not truncated)
            self._emit_first(env_idx, next_state, float(done))
    elif len(q) == self.n_step:
        self._emit_first(env_idx, next_state, 0.)
\end{lstlisting}

\subsection{Soft target updates and other details}

Instead of copying $\theta$ into $\theta^-$ every $C$ steps, the repository
uses \emph{Polyak averaging} $\theta^- \gets \tau \theta + (1 - \tau)
\theta^-$ after every update, with $\tau = 0.005$: $\theta^-$ is an
exponential average of the recent parameters, with a time constant of
$1 / \tau = 200$ updates. The rewards are multiplied by $0.1$ (so that the
values are of order $1$ to $10$), the gradients are clipped, and the
exploration rate $\varepsilon$ decreases linearly from $1$ to $0.02$ over
the first quarter of the training.

\section{Rainbow and beyond}
\label{sec:rainbow}

Rainbow (Hessel et al., 2018) combines six extensions of DQN: the three
above (double, dueling, $n$-step) and the three of this section. On the
Atari benchmark, each contributes; prioritised replay, $n$-step targets and
the distributional loss contribute most.

\subsection{Prioritised experience replay}

Uniform sampling spends most updates on transitions that are already well
predicted. Prioritised replay (Schaul et al., 2016) samples transition $i$
with probability
\begin{equation}
  P(i) = \frac{p_i^\alpha}{\sum_j p_j^\alpha}, \qquad p_i = \abs{\delta_i} +
  \epsilon,
\end{equation}
where $\delta_i$ is its last TD error and $\alpha \in [0, 1]$ interpolates
between uniform ($\alpha = 0$) and fully prioritised sampling. Sampling
from $P$ instead of the uniform distribution changes the expected
gradient; the bias is corrected by weighting the loss of transition $i$ by
\begin{equation}
  w_i = \left( \frac{1}{N\, P(i)} \right)^\beta,
\end{equation}
normalised by $\max_j w_j$. With $\beta = 1$ the weighted gradient is
unbiased: $\E_{i \sim P}[w_i g_i] = \sum_i P(i) \frac{1}{N P(i)} g_i =
\frac{1}{N} \sum_i g_i$. In practice $\beta$ grows from about $0.4$ to $1$
during training. In Pacman, deaths are rare and their TD errors large:
prioritisation would replay them far more often than uniform sampling.

\subsection{Distributional reinforcement learning}

Instead of the expected return $q(s, a)$, distributional RL learns the
whole distribution of the return $Z(s, a)$, which satisfies the
\emph{distributional Bellman equation}
\begin{equation}
  Z(s, a) \overset{D}{=} R + \gamma\, Z(S', A'),
\end{equation}
an equality of distributions, with $(R, S')$ drawn from $p(\cdot \mid s,
a)$ and $A' \sim \pi(\cdot \mid S')$. For policy evaluation, the
corresponding operator is a $\gamma$-contraction in the maximal
Wasserstein distance (Bellemare, Dabney and Munos, 2017). Two
parametrisations are common.
\begin{itemize}
  \item \textbf{C51} represents $Z(s, a)$ as a categorical distribution on
    $51$ fixed atoms $z_1 < \dots < z_{51}$; the target distribution of
    $r + \gamma Z(s', a^\star)$, whose atoms $r + \gamma z_j$ fall between
    the fixed atoms, is projected back onto them (each mass is split
    between the two neighbouring atoms in proportion to the distances), and
    the network minimises the cross-entropy to it.
  \item \textbf{QR-DQN} (Dabney et al., 2018) represents $Z(s, a)$ by $N$
    values $\theta_1, \dots, \theta_N$, meant to be its quantiles at levels
    $\tau_i = (2i - 1) / (2N)$, learnt with the quantile regression loss
    below.
\end{itemize}

\begin{proposition}{Quantile regression}{quantile}
  Let $Z$ be a real random variable with a continuous distribution function
  $F$, and $\rho_\tau(u) = u \left( \tau - \mathbb{1}[u < 0] \right)$ for
  $\tau \in (0, 1)$. Then $\theta \mapsto \E[\rho_\tau(Z - \theta)]$ is
  convex and minimised at the $\tau$-quantile $F^{-1}(\tau)$.
\end{proposition}

\begin{proof}
  $\rho_\tau$ is convex (it is piecewise linear with slopes $\tau - 1 < 0 <
  \tau$), hence so is the expectation. Its derivative with respect to
  $\theta$ is $\E[-\tau\, \mathbb{1}[Z > \theta] + (1 - \tau)\,
  \mathbb{1}[Z < \theta]] = -\tau (1 - F(\theta)) + (1 - \tau) F(\theta) =
  F(\theta) - \tau$, which vanishes at $F(\theta) = \tau$.
\end{proof}

The action values are recovered as means of the distributions, and the
agent still acts greedily with respect to them. Why does learning the
distribution help, if only its mean is used? The extra targets act as
auxiliary tasks that shape the network's representation, and they keep
apart situations with the same mean but different risks; in Pacman, a
move with a small chance of death and a safe move of equal mean value
have very different return distributions.

\subsection{Noisy networks}

$\varepsilon$-greedy exploration is the same in every state. NoisyNets
(Fortunato et al., 2018) replace the last linear layers by layers with
learnt noise, $y = (\mu^w + \sigma^w \odot \varepsilon^w) x + \mu^b +
\sigma^b \odot \varepsilon^b$, the noise $\varepsilon$ being resampled
regularly; the network learns how much to explore through $\sigma$, and
can explore differently in different states.

\subsection{Beyond Rainbow}

Recurrent networks (R2D2, Kapturowski et al., 2019) handle partial
observability, as in game \code{v2} (\cref{rem:pomdp}), by giving the
agent a memory; distributed actors (Ape-X) collect experience at a much
larger scale; Agent57 (2020) adds learnt exploration bonuses
(\cref{ch:exploration}) and surpassed the human benchmark on all $57$
Atari games. MuZero (\cref{ch:model-based}) replaces the replay of real
transitions by planning with a learnt model.

\section{The DQN agent of the repository}

\begin{pacman}{Network and training}{dqn-repo}
  \begin{itemize}[nosep]
    \item \textbf{Input:} the feature planes of the last two images
      ($14 \times 9 \times 13$, \cref{sec:observation}).
    \item \textbf{Network} (\code{DuelingCNN}): three $3 \times 3$
      convolutions with $32$, $64$ and $32$ channels, stride $1$ and
      padding $1$, a fully connected layer of $256$ units, and the two
      dueling heads; about $10^6$ parameters. The convolutions keep the full
      $9 \times 13$ resolution: on a board this small, a convolution of
      stride $2$ would blur the exact cells of Pacman and the ghosts, which
      decide everything.
    \item \textbf{Training:} $16$ games are played in parallel; at each step
      the network chooses the actions of the $16$ games in a single batched
      forward pass, and performs one gradient step on a mini-batch of $128$
      transitions from a memory of $300\,000$. Double DQN, $3$-step targets,
      Huber loss, $\tau = 0.005$, learning rate decreasing from $2.5 \times
      10^{-4}$ to $5 \times 10^{-5}$.
  \end{itemize}
\end{pacman}

\begin{lstlisting}[caption={One DQN update (\code{algorithms.py},
  \code{DQN.learn}, excerpt)}]
states, actions, returns, next_states, dones, discounts = \
    self.memory.sample(self.batch_size)
q = self.policy_net(states).gather(1, actions.unsqueeze(1)).squeeze(1)

with torch.no_grad():                      # targets held constant
    next_q_target = self.target_net(next_states)
    # double DQN: the online net chooses, the target net evaluates
    next_actions = self.policy_net(next_states).argmax(dim=1,
                                                       keepdim=True)
    next_q = next_q_target.gather(1, next_actions).squeeze(1)
    targets = (returns * self.reward_scale
               + discounts * (1. - dones) * next_q)  # n-step target

loss = F.smooth_l1_loss(q, targets)        # Huber loss
self.optimizer.zero_grad(set_to_none=True)
loss.backward()
nn.utils.clip_grad_norm_(self.policy_net.parameters(),
                         self.max_grad_norm)
self.optimizer.step()

with torch.no_grad():                      # Polyak averaging
    for p_t, p in zip(self.target_net.parameters(),
                      self.policy_net.parameters()):
        p_t.lerp_(p, self.tau)             # p_t <- p_t + tau (p - p_t)
\end{lstlisting}

\begin{remark}{The replay ratio}{replay-ratio}
  The \emph{replay ratio} is the average number of times each transition
  is used in a mini-batch. With the settings above, each step of the $16$
  games adds $16$ transitions and uses $128$: each transition is replayed
  about $8$ times. The first version of the repository played complete
  games in $32$ environments (up to about $32\,000$ transitions) and then
  made a \emph{single} gradient step on $32$ transitions: a replay ratio of
  about $0.001$. Its network barely learnt anything, whatever its
  architecture: in deep RL the amount of learning per transition matters as
  much as the algorithm.
\end{remark}

\begin{remark}{Results}{dqn-results}
  The results of the DQN agent trained for $40$ million steps on Kaggle
  will be added here, and compared with the A2C/PPO agent of
  \cref{ch:ppo}.
\end{remark}

\section{Summary}

\begin{itemize}[nosep]
  \item With function approximation, Q-learning becomes a semi-gradient
    method, which can diverge (deadly triad).
  \item DQN stabilises it with a replay memory (decorrelated, reused
    samples, possible because Q-learning is off-policy) and a slowly moving
    target network (each phase is a regression towards fixed targets).
  \item Double DQN removes most of the maximisation bias; the dueling
    architecture learns state values from every transition; $n$-step
    targets propagate rare rewards faster.
  \item In practice, the amount of learning per transition (replay ratio)
    is as important as these refinements.
\end{itemize}
